\section{问题四的建模与求解}
\subsection{固定任务下的独立分组问题}
问题四要求把15个服务区划为两个或三个独立救援组。这里以问题三已经验证的24条运输架次、5条中继架次及其全部开始和结束时刻为固定输入，不改变货箱组批、访问顺序、任务时刻或既定通信保障关系。每个服务区恰归属一个组，组间不能借用运输机、电池、中继机或能源组件。组内允许将同型号的实体设备重新编号，只要原任务时间表保持可执行。因而本问优化的对象是\emph{分组及各组需配置的设备数}，而不是重新求解第三问的联合调度。

独立执行有两层含义。其一，同一运输架次到达的所有服务区必须归入同一组，否则一条不可拆的既定路线会跨组。其二，若问题三的一条中继架次实际保障了若干运输路线，这些路线对应的服务区也必须归为同组；否则同一中继任务及其组件会被两个独立组同时使用。仅凭中继悬停点的几何覆盖范围合并服务区会过度约束分区，因此这里只采用\emph{连续通信证书中实际选用}的保障关系。

\subsection{不可拆分单元与可行分区}
以服务区为顶点构造无向图。若一条运输架次访问两个服务区，连接对应顶点；若某条中继任务在连续证书的至少一个已证区间实际保障运输架次$p$，则把$p$所访问的服务区与同一中继任务保障的其他架次的服务区连通。传递闭包给出不可拆单元。问题三证书的1317个区间无缝覆盖各运输架次的通信需求，其中547个区间采用中继保障，五条中继任务均有实际保障对象。长时中继任务先把九个服务区相连，后续中继任务经$S009$接入$S010,S013,S014$，再经$S014$接入$S011$；同架次跨区边补足联系。$S001$与$S006$均没有接入这些保障边。由此得到三个连通块：
\begin{equation}
\mathcal C_1=\{S001\},\qquad
\mathcal C_2=\{S006\},\qquad
\mathcal C_3=\{S002,\ldots,S015\}\setminus\{S006\}.
\label{eq:q4-components}
\end{equation}
任一独立组只能由这些连通块的并组成，且每组非空。三个无标号单元划成两组共有三种方式，划成三组只有一种；因此可逐一枚举，不需要启发式搜索。这里“无标号”表示仅交换组编号不会生成一个新方案。

\subsection{设备配置与评价准则}
固定第三问任务的时间轴后，同型号设备在组内可重编号，所需最少数量等于对应占用区间的最大并发数。对组$h$和资源种类$k$，记该组任务给出的半开区间集合为$\mathcal I_{hk}$，则
\begin{equation}
n_{hk}=\max_t\sum_{I\in\mathcal I_{hk}}\mathbf 1_I(t).
\label{eq:q4-peak}
\end{equation}
运输机的区间从准备开始到返航；同型共享电池延续到充满。中继机的区间从任务启动到返航后的300\,s周转结束；能源组件则从任务启动到返航充满。区间在右端点释放，故前一任务结束时可以衔接下一任务。分别对A、B、C三型运输机和电池、两类中继资源计算式~\eqref{eq:q4-peak}，没有把某架实体机的原编号误当作分组后的唯一可用机号。

设$N_k$为现有库存，组间配置总量$N_k^{\rm split}=\sum_h n_{hk}$。库存缺口、相对于问题三联合执行峰值$N_k^{\rm joint}$的新增配置量，分别定义为
\begin{equation}
\Delta_k=\max\{0,N_k^{\rm split}-N_k\},\qquad
R_k=N_k^{\rm split}-N_k^{\rm joint}.
\label{eq:q4-gap}
\end{equation}
这里$R_k$表示\emph{分组后相对于联合执行多占的设备数}，不是库存中闲置的设备。组$h$的运输作业量$W_h$取该组运输架次作业时长之和，用
\begin{equation}
I_K=\frac{\max_h W_h-\min_h W_h}{(\sum_h W_h)/K}
\label{eq:q4-balance}
\end{equation}
衡量$K$组之间的负担差异。作业量之和不是完工时间；同组并行作业仍逐项计入。对每种组数先最小化总缺口$\sum_k\Delta_k$，再最小化$I_K$，最后最小化$\sum_kR_k$。这一次序使“少补设备”优先于负担均衡，作为明确的决策偏好报告；若采用其他优先序，推荐分区可能变化。

\subsection{枚举结果与配置分析}
表~\ref{tab:q4-partition}列出按上述准则选择的分区。两组情形将$S001$单列，其余14区为另一组；三组情形只能把三个连通块分别成组。工作量较大的13区连通块承担全部运输作业时长的89.633\%，因此在不改动第三问通信保障关系时，不能把大部分任务均匀拆到不同组。
\begin{table}[htbp]
\centering\small
\caption{固定问题三任务下的独立分组与运输作业量}
\label{tab:q4-partition}
\begin{tabular}{@{}clrrrr@{}}
\toprule
组数 & 分组 & 货箱/箱 & 运输架次 & 中继架次 & 作业量/s\\
\midrule
2 & $\{S001\}$ & 15 & 2 & 0 & 3186.338\\
  & 其余14区 & 65 & 22 & 5 & 42611.904\\
\midrule
3 & $\{S001\}$ & 15 & 2 & 0 & 3186.338\\
  & 其余13区 & 59 & 21 & 5 & 41050.350\\
  & $\{S006\}$ & 6 & 1 & 0 & 1561.554\\
\bottomrule
\end{tabular}
\end{table}

两组方案中，单列$S001$需要1架C型机和2块C型电池；其余14区需要A/B/C型机$4/2/2$架、同型电池$6/4/3$块、中继机2架及组件3个。三组方案仅增加$S006$独立组所需的1架C型机和1块C型电池，13区主组所需的中继与A、B型资源均未减少。表~\ref{tab:q4-stock}把各组峰值相加后与现有库存比较。
\begin{table}[htbp]
\centering\small
\caption{独立配置总量与现有库存对照；斜线内依次为A/B/C型}
\label{tab:q4-stock}
\begin{tabular}{@{}lcccc@{}}
\toprule
方案 & 运输机/架 & 电池/块 & 中继机/架 & 组件/个\\
\midrule
现有库存 & $4/2/2$ & $6/4/4$ & 2 & 6\\
问题三联合执行峰值 & $4/2/2$ & $6/4/4$ & 2 & 3\\
分为两组所需 & $4/2/3$ & $6/4/5$ & 2 & 3\\
分为三组所需 & $4/2/4$ & $6/4/6$ & 2 & 3\\
\bottomrule
\end{tabular}
\end{table}

图~\ref{fig:q4-inventory}突出显示两种方案的C型资源缺口：两组最少需要补1架C型运输机和1块C型电池，共2件；三组分别需补2架和2块，共4件。A、B型设备与中继资源均未形成缺口，组件在两种方案中均只需3个。当前库存若不补齐缺口，就不能在\emph{保持第三问时间表不变}的条件下实现这些组的完全独立执行；这不是已可直接实施的分组排程。

枚举也揭示了配置与均衡之间的取舍：两组的三种无标号分区中，单列$S001$和单列$S006$均缺2件，但后者的不均衡指标为1.863615，高于前者的1.721707；将$S001$与$S006$并为一组可把不均衡降至1.585321，却会把缺口增加到4件。三组唯一分区的不均衡指标为2.586702。两组的缺口分布为两种方案缺2件、一种缺4件；三组只有一个方案且缺4件。这给出了本固定任务模型下的穷举最优性证明，而非对允许重新安排第三问任务时的所有分组策略给出下界。

\begin{figure}[htbp]
\centering
\includegraphics[width=0.93\textwidth]{figures/result_q4_inventory_comparison.pdf}
\caption{现有库存与两组、三组独立配置需求；背景标出形成缺口的C型运输机与电池}
\label{fig:q4-inventory}
\end{figure}

\subsection{适用范围与实施含义}
上述结论依赖于第三问正式方案、连续证书选定的实际中继保障边和既定任务时刻。若允许重排运输路线、改换中继位置或重新分配通信保障关系，连通块与资源峰值都可能改变，必须重新联合求解与验证。对于本题固定基线，首先补齐C型设备缺口，再按表~\ref{tab:q4-partition}划分，可在保留第三问已证安全、时限及连续通信性质的同时实现组间资源隔离。
